% !TEX program = lualatex
% Odosseus hackathon deck, LaTeX version of ../odysseus-presentation.html.
% Build:  latexmk          (uses ./latexmkrc: LuaLaTeX, output odosseus.pdf)
% Theme:  beamerthemeodosseus.sty  (colors, fonts and the slide building blocks)
\documentclass[aspectratio=169]{beamer}   % 160 x 90 mm = the HTML's 1600 x 900 px
\usetheme{odosseus}

\begin{document}

% ====================== SLIDE 1 · LOGO ======================
% The car image is © Constructor (from slide 2 of the challenge INTRO deck) and licensed for the hackathon only:
% it is kept out of git (see .gitignore); without the file the slide shows a short "not found" note instead.
\renewcommand\logocaption{"He is never at a loss"}   % caption, centred under the logo and the car
\renewcommand\logoside{images/a2rl_constructor_car_side.png}
\renewcommand\logocredits{Car image © Constructor 2026 (constructor.org), from the challenge INTRO slides, used for the National AI Hackathon 2026
  \quad·\quad The Odosseus logo was generated with ChatGPT (OpenAI)}
\logoslide

% ====================== SLIDE 2 · SOLUTION (final method, = ../../assets/approach.svg) ======================
\begin{odslide}
  \kicker{02}{Our solution}
  \slidetitle{Car state predicts, \highlight{the map corrects}}
  \slidesub{A causal filter: dead reckoning carries the pose between frames, camera fixes against the night map correct it.}

  % 5 steps + 4 arrows fill 1408 px: width factors sum to (1408 - 4 x 40) / 292.7 = 4.264
  \begin{pipeline}[210]     % [step height in px]
    \pipestep{plain}{Offline}{3D feature map}{SuperPoint + LightGlue, triangulated; MegaLoc database.}
    \pipearrow
    \pipestep[0.88]{plain}{Predict}{Automotive Dead reckoning}{Car-state speed and yaw rate carry the pose.}
    \pipearrow
    \pipestep[0.82]{hl}{Match}{3D keypoints}{Map points near the prediction, all 7 cameras.}
    \pipearrow
    \pipestep[0.7]{hl}{Solve}{PnP}{One robust PnP over all 7 cameras.}
    \pipearrow
    \pipestep[0.864]{new}{Update}{EKF}{No fix for 1 s: MegaLoc relocalises the pose.}
  \end{pipeline}

  \begin{stats}
    \stat{0.089}{m}{Position RMSE, held-out\\flying lap L7}
    \stat{0.30}{m}{Hardest lap L2,\\off the mapped lines}
    \stat{1–3.4}{s}{Recovery from 3–80 m\\injected pose errors}
    \stat{37}{ms}{Per 7-camera frame\\on one H100 NVL}
  \end{stats}
  \slidenote{\textbf{Relocalisation:} MegaLoc runs only after 1 s without an accepted camera fix; its pose is taken only if the jump is plausible for the time lost or has at least 300 inliers.}
\end{odslide}

% ====================== SLIDE 3 · ENGINEERING DECISIONS ======================
\begin{odslide}
  \kicker{03}{Engineering decisions}
  \slidetitle{What we dropped, and \highlight{why this works}}

  \begin{cards}[330]        % [card height in px]
    \card{no}{Didn't work}{
      \point{\textbf{Learned SLAM} (MASt3R-SLAM, VGGT-Long): lost track or 313 m error. 4 m between frames at night is too much.}
      \point{\textbf{Visual odometry}: best camera-only VO 43 m over 2 laps, vs 20.5 m from car state alone.}
      \point{\textbf{Track-border alignment}: cameras sit 0.4 m high, so 0.5 m sideways moves the border \textasciitilde 1 px.}
      \point{\textbf{Trusting odometry everywhere}: global down-weighting of fixes doubled the error on a well-mapped lap.}}
    \card{ok}{Why we chose this}{
      \point{\textbf{Car state for motion}: 4× more accurate than VO over 1–10 s. Cameras do what odometry can't: absolute fixes.}
      \point{\textbf{Prediction makes matching cheap}: guided search replaces retrieval, 410 \ra{} 41 ms.}
      \point{\textbf{All 7 cameras in one PnP}: with one camera, yaw and sideways shift blur together (0.5 \ra{} 0.12 m).}
      \point{\textbf{Down-weight only weak fixes}: keeps accuracy where the map is good, helps where it isn't.}}
  \end{cards}

  \insight{Key insight}{On mapped laps the floor is the ground truth, not the matcher: every parameter sweep lands at \textbf{0.103–0.108 m}, next to the GT's own \textasciitilde 0.1 m.}
\end{odslide}

% ====================== SLIDE 4 · ABLATIONS AND STRESS TESTS ======================
\begin{odslide}
  \kicker{04}{Ablation studies}
  \slidetitle{Ablations and \highlight{stress tests}}
  \begin{ablations}[210]{326}{4}    % [card height px]{card width px}{cards per row}
    \ablation{Odometry}{No visual odometry}
      {Car-state dead reckoning beats our best camera-only VO (43 m).}
      {20.5 m over 2 laps}
    \ablation{Cross-validation}{Plausible jumps only}
      {A jump must fit the time without a fix, or have \ensuremath{\geq} 300 inliers.}
      {F2/L2: 3.7 \ra{} 0.30 m}
    \ablation{Kidnap test}{Kidnap recovery}
      {A bad fix or a long blind stretch, injected into the running filter.}
      {3–80 m, 20°: 1.0–3.4 s}
    \ablation{Runtime}{Cached 3D points}
      {Triangulated once offline; each frame only matches the cached points.}
      {410 \ra{} 41 ms per frame}
  \end{ablations}

  \slidefigure{images/l2_position_lateral_error.png}{1200}
    {Hardest lap L2 (map without L2, Exp21 run): position error (top) and lateral error (bottom) over the lap. The 1.7 m peak is almost entirely lateral, where the car drives 8.5 m from every map lap.}
\end{odslide}

% ====================== SLIDE 5 · VISUALIZATION ======================
% Links: run:videos/... plays the local copy next to the PDF; replace with an https:// link once uploaded.
\begin{odslide}
  \kicker{05}{Visualization}
  \slidetitle{Visualization, \highlight{live at 10 Hz}}

  \fullvideo{images/viz_mapped_poster.jpg}{run:videos/mapped_L7_0-30s.mp4}{MAPPED MODE · NIGHT · CLICK TO PLAY}
\end{odslide}

% ====================== SLIDE 6 · SUMMARY (approach diagram = ../../assets/approach.svg, as images/approach.png) ======================
\summaryfigslide{06}{Thank you}{images/approach.png}{
  \sumstack{Predict,}{The car state carries the pose between frames.}
  \sumstack{Localize,}{Camera fixes against a map of the night laps correct it.}
  \sumstack{Tune}{On a held-out lap, checked on laps the map has never seen.}}

\end{document}
